% VI/43: Supplementary reader bridge. Original problems and proofs remain intact.
\paragraph{Why degree three is the first elliptic case.}
A plane cubic is the zero scheme of a homogeneous polynomial
$F(X,Y,Z)$ of degree three in $\mathbb P^2_k$.  Degree alone does
not guarantee smoothness or irreducibility.  For a \emph{smooth}
plane curve of degree $d$ the genus formula gives
\[
 g=\frac{(d-1)(d-2)}{2};\qquad d=3\ \Longrightarrow\ g=1.
\]
A smooth projective genus-one curve equipped with a chosen
$k$-rational point is an elliptic curve.  A singular cubic is
\emph{not} an elliptic curve even if its equation resembles a
Weierstrass equation.  Over a general field, the chosen rational
point is part of the definition and cannot be inferred from genus
alone.

\paragraph{Worked example 1: the point at infinity and smoothness.}
Assume $\operatorname{char}k\ne2,3$ and consider
\[
 E:\quad Y^2Z=X^3+aXZ^2+bZ^3,
 \qquad a,b\in k.
\]
On the chart $Z\ne0$ this becomes $y^2=x^3+ax+b$.
At infinity $Z=0$ forces $X=0$, leaving the unique point
$O=[0:1:0]$.  For
$F=Y^2Z-X^3-aXZ^2-bZ^3$ the partial derivative
$F_Z=Y^2-2aXZ-3bZ^2$ equals $1$ at $O$.  Hence this point
is smooth.  The cubic is smooth everywhere precisely when
\[
 4a^3+27b^2\ne0.
\]
Indeed a singular affine point would satisfy
$2y=0$, $3x^2+a=0$, and $x^3+ax+b=0$; eliminating $x$
gives the stated discriminant condition.  The nonsingularity
hypothesis is essential before introducing the elliptic group law.

\paragraph{Worked example 2: adding three explicit rational points.}
Over a field of characteristic $0$, take
\[
 E:\ y^2=x^3-x=x(x-1)(x+1).
\]
Its discriminant is nonzero and the points
$P=(0,0)$, $Q=(1,0)$, $R=(-1,0)$ all lie on $E$.
The line $y=0$ meets the cubic exactly at these three points,
counted with intersection multiplicity.  The chord-and-tangent
law says that collinear points satisfy $P+Q+R=O$.
For a Weierstrass equation the inverse is $(x,y)\mapsto(x,-y)$,
so $R=-R$.  Consequently
\[
 P+Q=R=(-1,0),\qquad 2P=2Q=2R=O.
\]
This example computes an actual sum; the three nonzero points
are the full rational $2$-torsion subgroup when the base field
contains $-1$, $0$, and $1$ (as it does here).

\paragraph{Worked example 3: tangent and doubling formulas.}
At a point $P=(x_1,y_1)$ with $y_1\ne0$, implicit
differentiation gives the tangent slope
\[
 m=\frac{3x_1^2+a}{2y_1}.
\]
The third intersection of the tangent with the cubic has
$x$-coordinate $m^2-2x_1$, and reflecting across the $x$-axis
yields $2P$.  Thus
\[
 x(2P)=m^2-2x_1,\qquad
 y(2P)=m\bigl(x_1-x(2P)\bigr)-y_1.
\]
For $y_1=0$ the tangent is vertical, and the correct conclusion
is $2P=O$ rather than division by zero.  For distinct points
with different $x$-coordinates one uses the chord slope
$m=(y_2-y_1)/(x_2-x_1)$ and
$x(P+Q)=m^2-x_1-x_2$.  The denominators mark separate
geometric cases that must be handled explicitly.

\paragraph{Worked example 4: a cuspidal cubic is not elliptic.}
Let $C\subset\mathbb P^2_k$ be the cubic
\[
 C:\quad Y^2Z=X^3.
\]
At $[0:0:1]$ all partial derivatives of $Y^2Z-X^3$ vanish,
so this is a singular point.  On the affine chart $Z=1$ the
parametrization $x=t^2$, $y=t^3$ shows how the cusp arises:
two coordinates are powers of one parameter, and the local
ring $k[t^2,t^3]$ is not integrally closed.  Its normalization
is $k[t]$.  The homogeneous map
\[
 [S:T]\longmapsto[ST^2:T^3:S^3]
\]
extends the parametrization to $\mathbb P^1$ and maps
$[0:1]$ to the point at infinity.  Although $C$ has arithmetic
genus one, its normalization has geometric genus zero.  The
elliptic group law of a smooth plane cubic therefore cannot be
transferred to the cusp without substantial modification.

\paragraph{Scheme-theoretic caution.}
The equations of a cubic, its singular points, and the tangent
cone belong to the scheme, not merely its visible drawing.
On the affine chart $Z=1$, smoothness of $F(x,y,1)=0$ at a
$k$-rational point means that at least one partial derivative
is nonzero there.  If both vanish, the Zariski tangent space has
dimension two rather than the expected dimension one.  This
connects the plane-cubic calculations back to VI/32.
